\chapter*{General Conclusion}
\label{chap:conclusion}

\addcontentsline{toc}{chapter}{Conclusion}

This thesis set out from a simple observation about the state of software
security practice. The community has learned to detect many vulnerability
classes with considerable precision, yet two hard problems remain
insufficiently addressed. The first is the detection of authorization flaws,
whose defining feature is often the absence of a required protection rather
than the presence of a dangerous code pattern. The second is the repair of
confirmed injection vulnerabilities, where an edit must be trusted not only
to silence a warning, but also to preserve the intended behaviour of the
program. The work reported here therefore studies both sides of the secure
software pipeline: how to make hard-to-see vulnerabilities explicit, and how
to repair a supported vulnerability only when the repair can be justified.

The first part of the document established the conceptual ground on which
these questions rest. We surveyed the landscape of software security, the
main families of vulnerability detection techniques, and the principles of
Automated Program Repair, separating the recurring stages of fault
localisation, patch generation and patch validation. We then examined the
methodology of systematic mapping studies and described how a
human-in-the-loop, protocol-driven pipeline can automate their operational
stages while preserving the methodological guarantees that give a secondary
study its authority. That review exposed a double gap: existing detectors
struggle with absence-based authorization errors that require architectural
reasoning, while learning-based and pattern-based repairers offer broad
coverage without a principled guarantee that an accepted edit is correct.

The first technical contribution of the thesis addresses the detection side
of this gap. We developed the Privilege Flow Graph (PFG), a domain-specific
graph representation for REST APIs that makes authorization architecture
explicit. Instead of treating source code as a flat sequence of tokens, the
PFG captures the structural relation between endpoints, guards,
intermediate logic and terminal data operations. On top of this
representation, the proposed detector combines a GATv2-based neural model
with structural validation rules, so that the system can both identify
candidate authorization vulnerabilities and explain the missing protection
that makes a path unsafe. The contribution was evaluated on a corpus of
75 real-world CVEs across Django and Spring Boot applications, and it showed
that authorization flaws are learnable from graph structure when extraction
coverage is adequate, whereas general-purpose tools such as Semgrep and
CodeQL achieve near-zero recall on the same class of vulnerabilities.

The second technical contribution answers the remediation side. We developed
\irsam{}, a remediation engine built on a single uncompromising premise: it
repairs an injection vulnerability only when it can prove the repair sound
against a precise and decidable criterion, and otherwise it abstains,
returning a machine-readable reason rather than an unverified edit. The
engine realises this stance through a seven-stage pipeline organised around
an Intent Abstract Model, a binding function that relocates untrusted data
from the grammar layer to the protocol channel, and a five-gate validator
that certifies each accepted patch. The soundness argument at the heart of
the design is mechanised in Lean~4, with the principal parameterisation and
audit results closed in the proved core. Safe-API knowledge is externalised
into declarative, schema-validated binder catalogs, so that extending
coverage is a data change rather than a change to the trusted code.

The empirical evaluation supports the central claims of both contributions.
For detection, the PFG achieved complete recall on real-world applications
where extraction was complete, and the failure analysis showed that missed
cases were primarily caused by extractor coverage rather than by the graph
formalism itself. For remediation, \irsam{} repaired every repairable case in
a controlled testbed of one hundred SQL injection instances with a
re-checkable certificate and a perfect five-gate pass rate. On the OWASP
Benchmark it returned a sound decision for all in-scope findings, never
certifying an unsound patch, and the abstention reasons it produced were
concrete and grounded in the code rather than opaque failures. Together,
these results support the thesis that useful security automation should be
both explicit about what it detects and disciplined about what it changes.

We have been candid throughout about the boundaries of the work. On the
detection side, the main limitation is extraction coverage: the PFG schema is
framework-agnostic, but real applications use framework-specific middleware,
routing conventions and dynamic patterns that require additional extractor
engineering. On the remediation side, the supported repair fragment is
narrower than the full space of injection sinks, the bundled detector limits
the volume of findings acted upon, and a small number of per-framework
wire-protocol lemmas in the Lean development remain to be discharged. None
of these is a soundness hole; each is a boundary of scope that maps to a
concrete next step. The most immediate extensions are to broaden framework
support for PFG extraction, register a host rewriter for the shell
interpreter, pair \irsam{} with production-grade detectors through its
detector-agnostic ingest layer, and complete the mechanised guarantee end to
end across every catalog.

The wider lesson of the thesis generalises beyond the specific defects it
addresses. Automated security tools become trustworthy when their internal
representation matches the structure of the vulnerability they target, and
when their outputs are accompanied by evidence strong enough to be checked.
For authorization flaws, this means making missing protections visible at
the architectural level. For injection repair, it means proving before
patching and abstaining audibly when no proof is within reach. It is on this
combined conviction, detection that explains and remediation that certifies,
that the contribution of this work rests.